\documentclass[aspectratio=169,11pt]{beamer}
\usepackage{beamerucad}
\usepackage{listings}
\definecolor{ckw}{HTML}{2563AB}\definecolor{cst}{HTML}{2E8B57}\definecolor{ccm}{HTML}{8A8F98}
\definecolor{outbg}{HTML}{F4F5F7}\definecolor{outrule}{HTML}{2E8B57}
\definecolor{errbg}{HTML}{FBEDEC}\definecolor{errrule}{HTML}{C0392B}
\lstdefinestyle{py}{language=Python,basicstyle=\ttfamily\scriptsize,
 keywordstyle=\color{ckw}\bfseries,stringstyle=\color{cst},commentstyle=\color{ccm}\itshape,
 showstringspaces=false,columns=fullflexible,keepspaces=true,
 backgroundcolor=\color{ucadband!8},frame=leftline,framerule=2.5pt,rulecolor=\color{ucadband},
 xleftmargin=8pt,aboveskip=3pt,belowskip=1pt,basewidth=0.5em}
\lstdefinestyle{out}{basicstyle=\ttfamily\scriptsize\color{black!82},
 backgroundcolor=\color{outbg},frame=leftline,framerule=2.5pt,rulecolor=\color{outrule},
 xleftmargin=8pt,aboveskip=2pt,belowskip=1pt,columns=fullflexible,keepspaces=true}
\lstdefinestyle{err}{basicstyle=\ttfamily\scriptsize\color{black!82},
 backgroundcolor=\color{errbg},frame=leftline,framerule=2.5pt,rulecolor=\color{errrule},
 xleftmargin=8pt,aboveskip=2pt,belowskip=1pt,columns=fullflexible,keepspaces=true}
\lstset{style=py}
\newcommand{\resultat}{{\scriptsize\textbf{R\'esultat}~$\rightarrow$}\par\vspace{1pt}}

\title{Programmation Python — Séance 2}
\subtitle{Structures de données}
\author{Dr. El Hadji Bassirou TOURÉ}
\institute{Département de Mathématiques et Informatique\\ Faculté des Sciences et Techniques\\ Université Cheikh Anta Diop de Dakar}
\date{}
\setucadfoot{Programmation Python — Séance 2}
\begin{document}

\begin{frame}[plain,noframenumbering]\titlepage\end{frame}

% =====================================================================
\begin{frame}{Objectifs de la séance}
\begin{ucaddef}[Contenu]
{\small Jusqu'ici, une variable contenait \emph{une} valeur. On apprend à en regrouper \textbf{plusieurs} dans une même structure. Le choix de la bonne structure est une décision centrale en programmation : chacune a ses forces.}
\end{ucaddef}
\begin{itemize}\small
  \item \textbf{Listes} : collections ordonnées et modifiables.
  \item \textbf{Tuples} : collections figées.
  \item \textbf{Dictionnaires} : associations clé $\rightarrow$ valeur.
  \item \textbf{Ensembles} : collections de valeurs uniques.
  \item \textbf{Compréhensions} : construire une liste en une ligne.
\end{itemize}
{\small L'entraînement se fait dans le \textbf{notebook}, avec une cellule « À vous » après chaque notion.}
\end{frame}

\begin{frame}{Panorama des quatre structures}
\figslide{0.92}{figs/f2_structures}{Quatre façons de regrouper des valeurs : par position (liste, tuple), par clé (dictionnaire), ou en collection unique (ensemble).}
\end{frame}

% #####################################################################
\section{Les listes}

\begin{frame}{L'idée : une collection ordonnée et modifiable}
\begin{ucaddef}[L'idée en une phrase]
{\small Une \textbf{liste} regroupe des valeurs entre crochets \texttt{[ ]}, dans un \textbf{ordre}. On y accède par \textbf{position}, comme pour une chaîne (indice à partir de \texttt{0}, \texttt{-1} pour le dernier), et — nouveauté — on peut la \textbf{modifier} après création.}
\end{ucaddef}
{\small Une liste peut contenir n'importe quoi : des nombres, des textes, et même un mélange. C'est la structure la plus utilisée pour stocker une suite de données du même genre — une liste de villes, de notes, de patients.}
\end{frame}

\begin{frame}[fragile]{Créer, accéder, découper}
\begin{lstlisting}
villes = ["Dakar", "Louga", "Mbour", "Kaolack"]
print(villes)
print("nombre :", len(villes))
print("premiere :", villes[0])
print("derniere :", villes[-1])
print("tranche  :", villes[1:3])
\end{lstlisting}
\resultat
\begin{lstlisting}[style=out]
['Dakar', 'Louga', 'Mbour', 'Kaolack']
nombre : 4
premiere : Dakar
derniere : Kaolack
tranche  : ['Louga', 'Mbour']
\end{lstlisting}
\end{frame}

\begin{frame}[fragile]{Modifier une liste}
\begin{ucadex}[Remplacer, ajouter, insérer]
{\small \texttt{liste[i] = ...} remplace ; \texttt{.append(x)} ajoute à la fin ; \texttt{.insert(i, x)} insère à la position \texttt{i}.}
\end{ucadex}
\begin{lstlisting}
villes[0] = "Thies"
villes.append("Touba")
villes.insert(1, "Saint-Louis")
print(villes)
\end{lstlisting}
\resultat
\begin{lstlisting}[style=out]
['Thies', 'Saint-Louis', 'Louga', 'Mbour', 'Kaolack', 'Touba']
\end{lstlisting}
\end{frame}

\begin{frame}[fragile]{Retirer un élément}
\begin{ucadex}[Par valeur ou par position]
{\small \texttt{.remove(x)} retire la valeur \texttt{x} ; \texttt{.pop()} retire le dernier élément \emph{et} le renvoie.}
\end{ucadex}
\begin{lstlisting}
villes.remove("Touba")
dernier = villes.pop()
print("retire :", dernier)
print(villes)
\end{lstlisting}
\resultat
\begin{lstlisting}[style=out]
retire : Kaolack
['Thies', 'Saint-Louis', 'Louga', 'Mbour']
\end{lstlisting}
\end{frame}

\begin{frame}[fragile]{Trier, compter, tester}
\begin{lstlisting}
nombres = [5, 2, 8, 1]
nombres.sort()
print(nombres)
print("taille :", len(nombres))
print(8 in nombres)
\end{lstlisting}
\resultat
\begin{lstlisting}[style=out]
[1, 2, 5, 8]
taille : 4
True
\end{lstlisting}
{\small \texttt{len()} compte les éléments ; \texttt{in} teste la présence d'une valeur et renvoie un booléen.}
\end{frame}

\begin{frame}[fragile]{Piège — \texttt{sort()} ne renvoie rien}
\begin{ucadpiege}[Modification en place]
{\small \texttt{.sort()} trie la liste \textbf{sur place} et renvoie \texttt{None}. Écrire \texttt{nombres = nombres.sort()} efface donc la liste !}
\end{ucadpiege}
\begin{lstlisting}
valeurs = [3, 1, 2]
resultat = valeurs.sort()
print("renvoye par sort() :", resultat)
print("liste triee        :", valeurs)
\end{lstlisting}
\resultat
\begin{lstlisting}[style=out]
renvoye par sort() : None
liste triee        : [1, 2, 3]
\end{lstlisting}
\end{frame}

% #####################################################################
\section{Les tuples}

\begin{frame}{L'idée : une collection figée}
\begin{ucaddef}[L'idée en une phrase]
{\small Un \textbf{tuple} ressemble à une liste — ordonné, accès par position — mais il est \textbf{figé} : une fois créé, on ne peut plus le modifier. On l'écrit entre parenthèses \texttt{( )}.}
\end{ucaddef}
{\small On l'emploie pour un groupe de valeurs qui forment un tout cohérent et ne doivent pas changer : des coordonnées \texttt{(latitude, longitude)}, une date \texttt{(jour, mois, annee)}. Le caractère figé est une \emph{garantie} : personne ne modifiera ces valeurs par erreur.}
\end{frame}

\begin{frame}[fragile]{Accéder à un tuple, et son piège}
\begin{lstlisting}
coord = (14.69, -17.44)
print(coord[0])
print(coord[1])
\end{lstlisting}
\resultat
\begin{lstlisting}[style=out]
14.69
-17.44
\end{lstlisting}
{\small Tenter de modifier un tuple provoque une erreur :}
\begin{lstlisting}[style=err]
coord[0] = 15.0
TypeError: 'tuple' object does not support item assignment
\end{lstlisting}
\end{frame}

\begin{frame}[fragile]{Déballer un tuple (unpacking)}
\begin{ucadex}[Affecter plusieurs variables d'un coup]
{\small On peut répartir les éléments d'un tuple dans autant de variables, en une ligne.}
\end{ucadex}
\begin{lstlisting}
latitude, longitude = coord
print("lat =", latitude)
print("lon =", longitude)
\end{lstlisting}
\resultat
\begin{lstlisting}[style=out]
lat = 14.69
lon = -17.44
\end{lstlisting}
{\small Très pratique pour échanger deux variables (\texttt{a, b = b, a}) ou récupérer plusieurs valeurs à la fois.}
\end{frame}

% #####################################################################
\section{Les dictionnaires}

\begin{frame}{L'idée : associer une clé à une valeur}
\begin{ucaddef}[L'idée en une phrase]
{\small Un \textbf{dictionnaire} associe des \textbf{clés} à des \textbf{valeurs}, entre accolades \texttt{\{ \}}. On n'accède plus par position, mais par \textbf{clé} : \texttt{patient["nom"]}. C'est la structure idéale pour décrire un objet par ses attributs.}
\end{ucaddef}
{\small Chaque clé est unique et pointe vers une valeur. Là où une liste répond à « quel est le troisième élément ? », un dictionnaire répond à « quelle est la valeur du nom ? » — bien plus lisible quand les données ont une étiquette naturelle.}
\end{frame}

\begin{frame}[fragile]{Créer et accéder}
\begin{lstlisting}
patient = {"nom": "Fatou", "age": 30, "ville": "Dakar"}
print(patient)
print("nom   :", patient["nom"])
print("ville :", patient["ville"])
\end{lstlisting}
\resultat
\begin{lstlisting}[style=out]
{'nom': 'Fatou', 'age': 30, 'ville': 'Dakar'}
nom   : Fatou
ville : Dakar
\end{lstlisting}
\end{frame}

\begin{frame}[fragile]{Ajouter et modifier}
\begin{lstlisting}
patient["tel"] = "77 000 00 00"
patient["age"] = 31
print(patient)
\end{lstlisting}
\resultat
\begin{lstlisting}[style=out]
{'nom': 'Fatou', 'age': 31, 'ville': 'Dakar', 'tel': '77 000 00 00'}
\end{lstlisting}
{\small Affecter une clé existante la \textbf{met à jour} ; une clé nouvelle est \textbf{créée}. Ici \texttt{age} passe à 31 et \texttt{tel} apparaît.}
\end{frame}

\begin{frame}[fragile]{Clé absente : \texttt{.get} et \texttt{in}}
\begin{ucadpiege}[Accès direct à une clé inexistante]
{\small \texttt{patient["email"]} sur une clé absente lève \texttt{KeyError: 'email'}. La méthode \texttt{.get(cle, defaut)} renvoie une valeur de repli au lieu de planter ; \texttt{in} teste la présence d'une clé.}
\end{ucadpiege}
\begin{lstlisting}
print(patient.get("email", "non renseigne"))
print("nom" in patient)
print("email" in patient)
\end{lstlisting}
\resultat
\begin{lstlisting}[style=out]
non renseigne
True
False
\end{lstlisting}
\end{frame}

\begin{frame}[fragile]{Parcourir clés, valeurs, couples}
\begin{lstlisting}
print(patient.keys())
print(patient.values())
print(patient.items())
\end{lstlisting}
\resultat
\begin{lstlisting}[style=out]
dict_keys(['nom', 'age', 'ville', 'tel'])
dict_values(['Fatou', 31, 'Dakar', '77 000 00 00'])
dict_items([('nom', 'Fatou'), ('age', 31), ('ville', 'Dakar'), ('tel', '77 000 00 00')])
\end{lstlisting}
{\small \texttt{.keys()}, \texttt{.values()}, \texttt{.items()} servent surtout à \emph{parcourir} un dictionnaire — ce qu'on fera avec les boucles en Séance 3.}
\end{frame}

% #####################################################################
\section{Les ensembles}

\begin{frame}{L'idée : des valeurs uniques}
\begin{ucaddef}[L'idée en une phrase]
{\small Un \textbf{ensemble} (\texttt{set}) est une collection \textbf{non ordonnée} et \textbf{sans doublons}, entre accolades \texttt{\{ \}}. Il sert à éliminer les répétitions et à tester très rapidement l'appartenance.}
\end{ucaddef}
{\small Comme l'ordre n'est pas garanti, on n'accède \emph{pas} par indice. Pour afficher le contenu de façon stable, on le trie avec \texttt{sorted()}, qui renvoie une liste ordonnée. Construire un ensemble à partir d'une liste est la façon la plus simple de supprimer les doublons.}
\end{frame}

\begin{frame}[fragile]{Unicité et appartenance}
\begin{lstlisting}
liste = ["Dakar", "Louga", "Dakar", "Mbour", "Louga"]
uniques = set(liste)
print("nombre d'uniques :", len(uniques))
print("contenu trie     :", sorted(uniques))
print("Dakar present ?  :", "Dakar" in uniques)
\end{lstlisting}
\resultat
\begin{lstlisting}[style=out]
nombre d'uniques : 3
contenu trie     : ['Dakar', 'Louga', 'Mbour']
Dakar present ?  : True
\end{lstlisting}
\end{frame}

\begin{frame}[fragile]{Union et intersection}
\begin{ucadex}[Combiner deux ensembles]
{\small \texttt{|} réunit (tout) ; \texttt{\&} garde les éléments \textbf{communs}.}
\end{ucadex}
\begin{lstlisting}
a = {"Dakar", "Thies", "Louga"}
b = {"Louga", "Mbour", "Dakar"}
print("union        :", sorted(a | b))
print("intersection :", sorted(a & b))
\end{lstlisting}
\resultat
\begin{lstlisting}[style=out]
union        : ['Dakar', 'Louga', 'Mbour', 'Thies']
intersection : ['Dakar', 'Louga']
\end{lstlisting}
\end{frame}

% #####################################################################
\section{Les compréhensions}

\begin{frame}{L'idée : construire une liste en une ligne}
\begin{ucaddef}[L'idée en une phrase]
{\small Une \textbf{compréhension de liste} crée une liste de façon concise à partir d'une autre collection, selon le modèle \texttt{[ expression for element in collection ]}.}
\end{ucaddef}
{\small On lit : « pour chaque \texttt{element} de la \texttt{collection}, calcule \texttt{expression} et range-la dans la nouvelle liste ». La boucle \texttt{for} complète est détaillée en Séance 3 ; ici, elle apparaît seulement dans cette écriture compacte, très courante en science des données.}
\end{frame}

\begin{frame}[fragile]{Transformer chaque élément}
\begin{lstlisting}
nombres = [1, 2, 3, 4]
doubles = [n * 2 for n in nombres]
print(doubles)
\end{lstlisting}
\resultat
\begin{lstlisting}[style=out]
[2, 4, 6, 8]
\end{lstlisting}
{\small Chaque élément de \texttt{nombres} est multiplié par 2 ; le résultat est une \textbf{nouvelle} liste, sans toucher à l'originale.}
\end{frame}

\begin{frame}[fragile]{Filtrer avec une condition}
\begin{ucadex}[Ajouter un \texttt{if} à la fin]
{\small Seuls les éléments qui vérifient la condition sont conservés.}
\end{ucadex}
\begin{lstlisting}
villes = ["Dakar", "Saint-Louis", "Mbour", "Ziguinchor", "Louga"]
longues = [v for v in villes if len(v) > 5]
print(longues)
\end{lstlisting}
\resultat
\begin{lstlisting}[style=out]
['Saint-Louis', 'Ziguinchor']
\end{lstlisting}
{\small On garde les villes dont le nom dépasse 5 caractères. Transformer et filtrer en une ligne rend le code court et lisible.}
\end{frame}

% #####################################################################
\section{Synthèse}

\begin{frame}{Quelle structure choisir ?}
\begin{center}
\small
\begin{tabular}{l c c c c}
\textbf{Structure} & \textbf{Syntaxe} & \textbf{Accès} & \textbf{Modifiable} & \textbf{Doublons}\\[2pt]
\hline\\[-6pt]
Liste & \texttt{[ ]} & par indice & oui & autorisés\\[2pt]
Tuple & \texttt{( )} & par indice & non (figé) & autorisés\\[2pt]
Dictionnaire & \texttt{\{c: v\}} & par clé & oui & clés uniques\\[2pt]
Ensemble & \texttt{\{ \}} / \texttt{set()} & sans indice & oui & interdits\\
\end{tabular}
\end{center}
\vspace{4pt}
{\small \textbf{Repère} : une suite ordonnée $\rightarrow$ liste ; un groupe figé $\rightarrow$ tuple ; des données étiquetées $\rightarrow$ dictionnaire ; des valeurs uniques $\rightarrow$ ensemble.}
\end{frame}

\begin{frame}{À retenir — Séance 2}
\begin{ucadret}[L'essentiel]
{\small
\textbf{Liste} \texttt{[ ]} : ordonnée, modifiable ; \texttt{.append}, \texttt{.insert}, \texttt{.remove}, \texttt{.pop}, \texttt{.sort} (en place, renvoie \texttt{None}), \texttt{len}, \texttt{in}.\\[2pt]
\textbf{Tuple} \texttt{( )} : figé ; accès par indice ; \emph{unpacking} \texttt{a, b = t}.\\[2pt]
\textbf{Dictionnaire} \texttt{\{ \}} : clé $\rightarrow$ valeur ; \texttt{d[cle]}, \texttt{.get(cle, defaut)}, \texttt{in}, \texttt{.keys/.values/.items}.\\[2pt]
\textbf{Ensemble} \texttt{set()} : valeurs uniques, non ordonnées ; \texttt{|} union, \texttt{\&} intersection ; \texttt{sorted()} pour afficher.\\[2pt]
\textbf{Compréhension} : \texttt{[expression for element in collection if condition]}.}
\end{ucadret}
\end{frame}

\begin{frame}{Pièges fréquents}
\begin{itemize}\small
  \item \texttt{nombres = nombres.sort()} : \texttt{.sort()} renvoie \texttt{None} et efface la liste.
  \item Modifier un \textbf{tuple} : impossible (\texttt{TypeError}) — il est figé.
  \item Accéder à une clé absente d'un dictionnaire : \texttt{KeyError} ; préférer \texttt{.get}.
  \item Attendre un \textbf{ordre} dans un ensemble : il n'y en a pas ; trier avec \texttt{sorted()}.
  \item Confondre les délimiteurs : \texttt{[ ]} liste, \texttt{( )} tuple, \texttt{\{ \}} dictionnaire ou ensemble.
\end{itemize}
\end{frame}

\begin{frame}{Prochaine séance}
\begin{ucaddef}[Séance 3 — Contrôle du flux]
{\small On sait stocker des collections. La Séance 3 apprend à les \textbf{parcourir} et à \textbf{décider} : les conditions \texttt{if}/\texttt{elif}/\texttt{else}, et les boucles \texttt{for} et \texttt{while}. C'est là que la compréhension de liste révélera toute sa logique, et qu'on pourra traiter une liste élément par élément.}
\end{ucaddef}
{\small À faire d'ici là : refaire les cellules « À vous » du notebook de la séance.}
\end{frame}

\begin{frame}{Références}
\begin{itemize}\small
  \item Severance, C. — \emph{Python for Everybody}, 2016.
  \item Downey, A. — \emph{Think Python}, 2e éd., O'Reilly, 2015.
  \item Matthes, E. — \emph{Python Crash Course}, 3e éd., No Starch Press, 2023.
  \item Python Software Foundation — \emph{The Python Tutorial} (Data Structures), 2024.
\end{itemize}
\end{frame}

\end{document}
